\origpage{102}{ROBERTSON AND SEYMOUR}
Let $\operatorname{ins}'$ be the restriction of $\operatorname{ins}$ to the set of circuits of $\Gamma_i$ of length $<2\theta$. This is a slope of order $\theta$ in $\Gamma_i$, by (4.2). Let us say that $X$ is an $\operatorname{ins}'$-restraint if it is a restraint of $\Gamma_i$ with respect to $\operatorname{ins}'$. Then $X$ is an $\operatorname{ins}'$-restraint if and only if $X$ is a restraint of $\Gamma$ with respect to $\operatorname{ins}$ and $bd(X) \subseteq U(\Gamma_i)$. But, by the argument above when $\Gamma$ is connected, the union of $\operatorname{ins}'$-restraints $X$ with $l(X) \le k$ and $c \in X$ is simply-connected. This union equals $Z_i$, and $Z_i = Z$, and so $Z$ is simply-connected. $\square$

\stmt
(8.11) \textit{Let $c \in \Sigma$, let $0 \le k < 2\theta$, and let $Z$ be the union of all restraints $X$ with $l(X) \le k$ and $c \in X$. Then for any restraint $Y$ with $c \in Y$, if $bd(Y) \subseteq Z$ then $Y \subseteq Z$.}

\textit{Proof.} Suppose that $Y$ is a restraint with $c \in Y$, $bd(Y) \subseteq Z$ and $Y \nsubseteq Z$. Since $\Sigma - Z$ is arc-wise connected, by (7.6) and (8.10), and $bd(Y) \cap (\Sigma - Z) = \emptyset$, it follows that $\Sigma - Z \subseteq Y$. Let $X$ be a restraint with $Y \subseteq X$ and $l(X) \le l(Y)$, chosen with $X$ maximal; then $X$ is extreme, $c \in X$, and $\Sigma - Z \subseteq X$ and hence $bd(X) \subseteq Z$. Let $Z'$ be the union of all restraints $X'$ with $l(X') \le k$, $c \in X'$ and $bd(X') \subseteq X$.

Suppose that $Z' \nsubseteq X$. Then there is a restraint $X'$ with $l(X') \le k$, $c \in X'$, $bd(X') \subseteq X$, and $X' \nsubseteq X$. It follows, since $\Sigma - X$ is arc-wise connected by (7.6), that $X \cup X' = \Sigma$, contrary to (8.4). Hence $Z' \subseteq X$, and we claim that $Z' = Z \cap X$. Since certainly $Z' \subseteq Z$, it follows that $Z' \subseteq Z \cap X$. On the other hand, let $z \in Z \cap X$, and let $X'$ be a $c, z$-optimal restraint. Then $l(X') \le k$. Since $c, z \in X$ it follows from (8.9) that $X' \subseteq X$, and so $bd(X') \subseteq X$. Thus $X' \subseteq Z'$, and so $z \in Z'$. We deduce that $Z \cap X \subseteq Z'$, and so $Z' = Z \cap X$.

Now $\Sigma - X$, $\Sigma - Z$ are non-empty open sets, and they are disjoint, since $\Sigma - Z \subseteq X$. Hence $(\Sigma - X) \cup (\Sigma - Z)$ is not arc-wise connected. By (7.6), $X \cap Z = Z'$ is not simply-connected. But this contradicts (8.10) applied to $\Gamma \cap X$. We deduce that there is no such $Y$, as required. $\square$

\stmt
(8.12) \textit{If $c \in \Sigma$, there is a region $r$ of $\Gamma$ such that there is no restraint $X$ with $c \in X$ and $r \subseteq X$.}

\textit{Proof.} The union of all restraints $X$ with $l(X) = 0$ and $c \in X$ is $\{c\}$ if $c \in V(\Gamma)$ and $\emptyset$ otherwise. Thus we may choose $k \ge 0$ maximum such that $k < 2\theta$ and $Z_k \ne \Sigma$, where $Z_k$ is defined to be the union of all restraints $X$ with $l(X) \le k$ and $c \in X$. Since $Z_k$ is closed, there is a region $r \subseteq \Sigma - Z_k$, and we claim that $r$ satisfies the theorem. For suppose not. Then $k \le 2\theta - 2$ and $Z_{k+1} = \Sigma$, and in particular there is a restraint $X$ with $l(X) \le k+1$ and $c \in X$ and $r \subseteq X$. Choose such a restraint, extreme. Since $X \nsubseteq Z_k$ it follows from (8.11) that $bd(X) \nsubseteq Z_k$. Let $b \in bd(X) - Z_k$. Since $b \in bd(X)$ and $X$ is closed, there is a region $r'$ of $\Gamma$ with $r' \subseteq \Sigma - X$ and $b \in \bar{r}'$. Since $b \notin Z_k$ it follows that $r' \cap Z_k = \emptyset$. Since $Z_{k+1} = \Sigma$, there is a $c, r'$-optimal restraint $X'$, and $l(X') = k+1$. Since $b \notin Z_k$ and $b \in X'$ it follows that $X'$ is
